\documentclass[10pt,a4paper]{article}
\usepackage{amsmath,amssymb,amsthm,graphicx}
\usepackage{amsmath}
\usepackage{enumerate}
\usepackage{amsfonts}
\usepackage{algpseudocode}
\usepackage{algorithm}
\usepackage{amsthm}
\newtheorem{theorem}{Theorem}
\newtheorem{proposition}[theorem]{Proposition}
\newtheorem{lemma}[theorem]{Lemma}
\newtheorem{corollary}[theorem]{Corollary}
\newtheorem{definition}[theorem]{Definition}
\newtheorem{example}[theorem]{Example}
\newcommand{\out}{\textup{out}} 
\newcommand{\spa}{\textup{span}} 
\newcommand{\Stab}{\textup{Stab}} 
\newcommand{\peb}{\textup{peb}} 

\newcommand{\AcceptEdge}[1]{\textbf{Accept-Edge}(#1)}
\newcommand{\MovePebble}[1]{\textbf{Move-Pebble}(#1)}
\title{The admissible edge sets of the next layer form a matroid that does not depend on the chosen maximally sparse set}
\author{Signe Lundqvist, Tovohery Randrianarisoa, Klara Stokes, Joannes Vermant}
\date{arxiv:2306.05050v3; CC BY 4.0}
\begin{document}
\maketitle
\noindent\textit{Editorial note.} The surplus-independence lemma at source lines 676-726 is the article's own source-local support result and is delivered with its complete proof; it is already prepared in the corpus as problem 142659 and is not a newly counted target here.\par\medskip
\noindent\textit{Editorial note.} The two counting lemmas invoked inside that support proof are supplied by their statements only (source lines 600-611 and 631-640), exactly as in the prepared item 142659; their proofs belong to that other unit and are not part of this task.\par\medskip
\noindent\textit{Editorial note.} The context span at source lines 564-578 begins with the notation sentence that defines V\_i(F); it continues a preceding line of the article that also points forward to the matroid theorem, which is not part of this task.\par\medskip
\noindent\textit{Editorial note.} No retained line contains a citation, a figure or an included file, so the package declares neither a bibliography nor an asset.\par\medskip
\noindent\textit{Editorial note.} The retained environments print with this package's own numbering rather than the numbering of the article's full document.\par\medskip
\noindent\textit{Editorial scope.} One original exact result: the matroid well-definedness of the admissible edge family of a layer, with the article's own complete surplus-independence lemma as source-local support.\par\medskip
\noindent\textit{Allowed standard prerequisites (editorial).}\begin{itemize}
\item Matroid independence axioms. A hereditary family with the augmentation property, equivalently a family determined by its circuits; no rank formula is assumed.
\item Maximally sparse sets and the surplus. Maximally sparse edge sets on V\_1 union ... union V\_j and the minimum-defining surplus functional theta\_\{alpha\_j\} are supplied by the retained context spans (source lines 399-406, 564-578); membership in I\_j is the criterion displayed at 575-577, not the naive edge count.
\end{itemize}\medskip
\begin{definition}\label{KL-sparse}
    Let $G=(V,E)$ be a (multi)graph such that $V$ has a partition into subsets $V_1, V_2, \dots, V_n$, where every edge $e\in E$ has endpoints in $V_i$  and $V_{i+1}$ for some $i\in \{1,\dots, n-1\}$.
    
       Let $K=(k_1,\dots, k_n)\in \mathbb{R}_{>0}^{n}$ and ${L=(l_{1,2}, \dots , l_{n-1, n}) \in \mathbb{R}_{\geq 0}^{n-1}}$. For any $V'\subseteq V$, we define
    \begin{equation*}
        l(V') = \underset{\substack{
       i: V_{i}\cap V'\neq \emptyset\\ 
       V_{i+1}\cap V'\neq \emptyset}}
       {\sup} \{l_{i,i+1}\}.
    \end{equation*}
    We say that $G$ is $(K,L)$-sparse if, for any subset $V' \subseteq V$, one has
    \begin{equation*}
       \vert E(V')\vert \leq \sum_{i=1}^{n} k_i\vert V' \cap V_i\vert -l(V').
    \end{equation*}
    We note that the use of $\sup$ makes $l(V')=-\infty$ whenever the defining set is empty.\\ 
    Furthermore, $G$ is $(K,L)$-tight if $G$ is $(K,L)$-sparse, and additionally
        \begin{equation*}
       \vert E\vert = \sum_{i=1}^{n} k_i\vert V \cap V_i\vert -\sup_{i=1,\dots, n-1} \{l_{i,i+1}\}.
    \end{equation*}
    We furthermore say that a subset of edges $E'\subseteq E$ is $(K, L)$-sparse or $(K, L)$-tight if the graph $(V, E')$ is $(K,L)$-sparse or $(K,L)$-tight, respectively.  We say that a subset of edges $E'\subseteq E$ is maximally sparse if $E'$ is $(K,L)$-sparse, and if for any edge $vw \in E \setminus E'$ such that $v,w \in V(E')$, $E'\cup \{vw\}$ is not $(K,L)$-sparse.
\end{definition}

We will mainly consider the case where the entries of $K$ and $L$ are integers, and we need the following inequalities to hold for the $(K,L)$-sparse graphs to be a family of graphs that is well-behaved. 
\begin{align}
    k_i + k_{i+1} &> l_{i, i+1} \label{conditionposet1 -k k l},\\
    k_{i+1} + l_{i-1, i} &> l_{i, i+1}\label{conditionposet2 - k l l},\\
     l_{1,2} \leq  \dots &\leq l_{n-1,n}. \label{conditionposet3-increasing}
\end{align}

As stated in the introduction, we will show that one does, however, get a greedoid. We will consider strings $e_1\dots e_m$ on the alphabet $E$. For any string $\alpha=e_1\dots e_m$, we define the edge set $E_\alpha:= \{e_1, \dots, e_m\}$. Then, we define the collection $\mathcal{L}_{(K,L)}$ to be the set of strings on $E$ that satisfy the following three properties:
\begin{enumerate}
\item{For any $\alpha=e_1\dots e_m$, the set $E_\alpha$ is $(K,L)$-sparse}.
\item{For any $\alpha=e_1\dots e_m$ and $1\leq i< j \leq n-1$, if $e_{k_1} \in E(V_i \cup V_{i+1})$ and $e_{k_2} \in  E(V_j \cup V_{j+1})$, then $k_1 < k_2$.}
\item{For any $\alpha=e_1\dots e_m$ and $i\in \{1,\dots, n-1\}$ such that there is an $e_j$ with $e_{j}\in E(V_i\cup V_{i+1})$, the subgraph $E_{\alpha} \cap E(V_1\cup \dots \cup V_i)$ is maximally $(K,L)$-sparse on $V_1\cup \dots \cup V_i$. }
\end{enumerate}

We will use the terminology \textit{maximal feasible set} to mean a set $E_{\alpha}$, where $\alpha \in \mathcal{L}_{(K,L)}$ is a maximal word, meaning there is no word $\beta \in \mathcal{L}_{(K,L)}$ such that $\beta = \alpha e$ for some $e$. Expanding on what this means, the maximal feasible sets $E'$ are the sets $E'\subseteq E$ such that $E'(V_1 \cup \dots \cup V_j)$ is a maximally sparse set for each $j\in \{2, \dots, n\}$.

To prove that $\mathcal{L}_{(K,L)}$ is a greedoid, we will show that the sets of edges that can be added to a word $\alpha$ between $V_j$ and $V_{j+1}$ are the independent sets of a matroid. Let us make this precise. Suppose that $E^*$ is a maximally sparse set supported on $V_1 \cup...\cup V_{j}$, and let $\mathcal{I}_j$ be the collection of subsets $E' \subseteq E(V_{j},V_{j+1})$ such that $E^* \cup E'$ is $(K,L)$-sparse. We will show that the elements of $\mathcal{I}_j$ are the independent sets of a matroid with ground set $E(V_{j},V_{j+1})$. We will moreover show (see Lemma \ref{Matroid_1}) that $\mathcal{I}_j$ does not depend on the chosen maximally sparse set. We denote the pair $(E(V_{j},V_{j+1}), \mathcal{I}_j)$ by $\mathcal{M}_j.$

for any subset of edges $F\subseteq E$, we define $V_i(F):=V(F)\cap V_i$.
Let $\alpha_j$ be a word in $\mathcal{L}_{(K,L)}$ such that no $e\in E(V_{j+1} \cup V_{j+2})$ appears in $\alpha_j$, and $E_{\alpha_j}$ is maximally sparse on $V_1 \cup \dots \cup V_{j+1}$. We then define, for subsets $X\subseteq V_{j+1}$ the \textit{surplus}
\begin{align*}
    \theta_{\alpha_j}(X) := \min_{W: X\subseteq W}\{(\sum_{i=1}^{j+1} k_{i}|V_{i} \cap W| - |E_{\alpha_j}(W)|\}.
\end{align*} 
We remark that $\theta$ depends on the chosen parameters $K$.We will sometimes suppress the subscript for $\alpha_j$, writing $\alpha$ in stead of $\alpha_j.$ It follows easily from the definition that if $X\subseteq Y$ then $\theta_{\alpha}(X) \leq \theta_{\alpha}(Y)$. 


Let us briefly motivate the definition of the surplus.
Fix some $\alpha \in \mathcal{L}_{(K,L)}$ such that $E_{\alpha} \subseteq E(V_1 \cup \dots \cup V_j)$ is a maximally sparse subset of $ E(V_1\cup \dots \cup V_j)$. Defining $\mathcal{I}_j$ using $E_{\alpha}$, one can see that a set $E'\subseteq E(V_j\cup V_{j+1})$ belongs to $\mathcal{I}_j$ if and only if
\begin{equation}
    |F|\leq k_{j+1}|V_{j+1}(F)| + \theta_{\alpha}(V_{j}(F)) - l_{j,j+1}
    \label{def:matroid}
\end{equation}
for any subset $F \subseteq E'$.  

\begin{lemma}\label{Increase_when_tight - greedoids}
Let $K\in \mathbb{Z}_{>0}^{n}, L\in \mathbb{Z}_{\geq 0}^{n-1}$ be such that inequalities \eqref{conditionposet1 -k k l}, \eqref{conditionposet2 - k l l} and \eqref{conditionposet3-increasing} are valid. Let $G=(V,E)$ be a graph such that the vertex set has a partition $V=V_1\cup \dots \cup V_n$ such that all edges $e\in E$ have endpoints in $V_{i} \cup V_{i+1}$ for some $1 \leq i < n$. 

Suppose also that $B \subseteq E$ is $(K,L)$-sparse. Let $X\subseteq V_{1}\cup \dots \cup V_{j}$ and $T \subseteq V_{1}\cup \dots \cup V_{j}$ be such that $X\cap T$ contains some vertices of $V_{j-1}$ and $V_{j}$. Suppose also that
\begin{equation*}
\sum_{i=1}^{j} k_i |V_i\cap T| - |B(T)|  = \ell_{j-1, j}. 
\end{equation*}
Then
\begin{equation*}
\sum_{i=1}^{j} k_i |V_i \cap X| - |B(X)|  \geq \sum_{i=1}^{j} k_i |V_i\cap (T\cup X)| - |B(T\cup X)|. 
\end{equation*}
\end{lemma}

\begin{lemma}\label{blocks}
Let $K\in \mathbb{Z}_{>0}^{n}, L\in \mathbb{Z}_{\geq 0}^{n-1}$ be such that inequalities \eqref{conditionposet1 -k k l}, \eqref{conditionposet2 - k l l} and \eqref{conditionposet3-increasing} are valid. Let $G=(V,E)$ be a graph such that the vertex set has a partition $V=V_1\cup \dots \cup V_n$ such that all edges $e\in E$ have endpoints in $V_{i} \cup V_{i+1}$ for some $1 \leq i < n$, and assume that $G$ is $(K,L)$-sparse. 

Let $B_1, B_2 \subseteq E$ be tight edge sets, with $V(B_i)\subseteq V_1\cup \dots \cup V_t$. If $V(B_1)\cap V(B_2) \cap V_{t-1} \neq \emptyset$ and $V(B_1)\cap V(B_2) \cap V_{t} \neq \emptyset$, then $B_1 \cap B_2$ and $B_1\cup B_2$ are tight subgraphs.
Moreover, one has
\begin{align*}
    V(B_1\cup B_2) &= V(B_1) \cup V(B_2)\\
    V(B_1\cap B_2) &= V(B_1) \cap V(B_2)
\end{align*}
\end{lemma}

\begin{lemma}\label{Matroid_1}
Let $K\in \mathbb{Z}_{>0}^{n}, L\in \mathbb{Z}_{\geq 0}^{n-1}$ be such that inequalities \eqref{conditionposet1 -k k l}, \eqref{conditionposet2 - k l l} and \eqref{conditionposet3-increasing} are valid. Let $G=(V,E)$ be a graph such that the vertex set has a partition $V=V_1\cup \dots \cup V_n$ such that all edges $e\in E$ have endpoints in $V_{i} \cup V_{i+1}$ for some $1 \leq i < n$.

Suppose that $\mathcal{M}_c$ is a matroid for all $c < j$. Then the set $\mathcal{I}_{j}$ is independent of the maximally sparse set $E^*$. Furthermore, $\mathcal{I}_j$ is non-empty whenever $E(V_{j} \cup V_{j+1})$ is non-empty.
\end{lemma}

\begin{lemma}
    Let $K\in \mathbb{Z}_{>0}^{n}, L\in \mathbb{Z}_{\geq 0}^{n-1}$ be such that inequalities \eqref{conditionposet1 -k k l}, \eqref{conditionposet2 - k l l} and \eqref{conditionposet3-increasing} are valid. Let $G=(V,E)$ be a graph such that the vertex set has a partition $V=V_1\cup \dots \cup V_n$ such that all edges $e\in E$ have endpoints in $V_{i} \cup V_{i+1}$ for some $1 \leq i < n$. 
    
    If $\mathcal{M}_c$ is a matroid for all $1 \leq c\leq j$, then $\theta_{\alpha_j} = \theta_{\alpha_j'}$ for every two words $\alpha_j$ and $\alpha_j'$, such that $E_{\alpha_j}$ and $E_{\alpha_j'}$ are maximal feasible subsets supported on $V_1\cup \dots \cup V_{j+1}$. 
    \label{surplus-independent}
\end{lemma}
\begin{proof}
We are given $\alpha_j$ and $\alpha_j'$. For any word $\alpha_j = e_1, \dots e_m$ such that $E_{\alpha_j}$ is supported on $V_1\cup \dots \cup V_{j+1}$ and $d<j$, define $\alpha_{d}$ to be the word $e_1 \dots e_{\tilde{m}}$ obtained by restricting to those edges supported on $V_1\cup \dots \cup V_{d+1}$. I.e. $\alpha_{d}$ is the word resulting from deleting the edges $e_{\tilde{m}+1}\dots e_{m}$ where $\tilde{m}+1$ is the smallest index such that for all $i\geq \tilde{m}+1$, one has $e_{i} \in E(V_{d+i}\cup V_{d+i+1})$ for some $i \geq 1$. In particular, $\alpha_{0}$ will always be the empty word and for subsets $X\subseteq V_1$, that $\theta_{\alpha_0}(X) = \theta_{\alpha_0'}(X) = k_1|X|$, which serves as the basis for our induction. Suppose now that the statement holds for all $d < j$. We will show it holds for $d=j$.

Let $X\subseteq V_{j+1}$. One can verify that one can express $\theta_{\alpha_{j}}$ as a function of $\theta_{\alpha_{j-1}}$. Namely, one has:
\begin{equation*}
    \theta_{\alpha_j}(X)= \min_{\substack{W\subseteq V_{j} \cup V_{j+1}\\[2pt] X\subseteq W}}  k_{j+1}|V_{j+1} \cap W| +\theta_{\alpha_{j-1}}(V_{j}\cap W)  -|E_{\alpha_j}(W)|.
\end{equation*}


By induction, we can assume that $\theta_{\alpha_{j-1}} = \theta_{\alpha_{j-1}'}$. Thus, it suffices to consider the case where $\alpha_{j}$ and $\alpha_{j}'$ only differ for edges $e\in E(V_{j}\cup V_{j+1})$, and we may assume $\alpha_{j}\neq \alpha_{j}'$ since otherwise the lemma is trivially true. In this case, since we know that $\mathcal{M}_{j}$ is a matroid, we only need to show that in the case where $E_{\alpha_{j}} = \left (E_{\alpha_{j}'} \cup \{e\}\right) \setminus \{f\}$, one has $\theta_{\alpha_{j}}(X) = \theta_{\alpha_{j}'}(X)$, since the statement then follows for all maximally chosen words $\alpha_j$ and $\alpha_j'$. Suppose that the added edge, $e$, has endpoints $v,w$. By Lemma \ref{blocks}, we know that the intersection 
\begin{equation*}
M  = \bigcap_{\substack{T\subseteq E_{\alpha_j'},\\ v,w \in V(T),\\ T \textup{ tight }}} T
\end{equation*}
is a nonempty tight subset of $E_{\alpha_j'}$.  Furthermore, if $f$ were not in $M$, then $M$ would be a tight subset of $E_{\alpha_j}$, as $f$ is the only edge of $E_{\alpha_j'}$ not in $E_{\alpha_j}$. Then, $M \cup \{e\}$ would not be $(K,L)$-sparse, which contradicts $E_{\alpha_j}$ being $(K,L)$-sparse. It follows that $f \in M$.  Now, let $X\subseteq V$, and suppose $W$ minimises the expression for $\theta_{\alpha_j'}$, so
\begin{equation*}
\theta_{\alpha_{j}'}(X) = \sum_{i=1}^{j+1} k_i|V_i\cap W| -|E_{\alpha_j'}(W)|.
\end{equation*}

 If $f \notin E_{\alpha_j'}(W)$, then we see that

 \[|E_{\alpha_j}(W)| \geq |E_{\alpha_j'}(W)|\]

so it follows immediately that 

\begin{align*}
    \theta_{\alpha_j}(X) &\leq  \sum_{i=1}^{j+1} k_i |V_i \cap W| - |E_{\alpha_{j}}(W)| \\
    &\leq \sum_{i=1}^{j+1} k_i |V_i \cap W| - |E_{\alpha_{j}'}(W)| = \theta_{\alpha_j'}(X).
\end{align*}
 
 If $f \in E_{\alpha_j'}(W)$, then we get using Lemma \ref{Increase_when_tight - greedoids}, that
\begin{align*}
\sum_{i=1}^{j+1} k_i |V_i\cap W| - |E_{\alpha_j'}(W)|  \geq \sum_{i=1}^{j+1} k_i |V_i\cap (W\cup V(M))| - |E_{\alpha_{j}'}(W \cup V(M))|,
\end{align*}
and hence we may assume that $V(M) \subseteq W$.

Since $E_{\alpha_{j}} = \left (E_{\alpha_{j}'} \cup \{e\}\right) \setminus \{f\}$, and $e$ and $f$ are both supported on $V(M)$ it follows that,
$$|E_{\alpha_{j}'}(W\cup V(M))|=|E_{\alpha_{j}}(W\cup V(M))|.$$

It now follows that

$$\theta_{\alpha_j}(X) \leq  \sum_{i=1}^{j+1} k_i |V_i \cap W| - |E_{\alpha_{j}'}(W)| = \theta_{\alpha_j'}(X).$$


We conclude that $\theta_{\alpha_j}(X)\leq \theta_{\alpha_j'}(X)$ in all cases. Switching the roles of $E_{\alpha_j}$ and $E_{\alpha_j'}$, the same argument gives that $\theta_{\alpha_j}(X) \geq \theta_{\alpha_j'}(X)$, so $\theta_{\alpha_j'}(X) = \theta_{\alpha_j}(X)$. 
\end{proof}

\begin{proof}

First, assuming that $E^* \cap E(V_{j-1}\cup V_j)$ is non-empty, it is easy to see, using inequality \eqref{conditionposet2 - k l l}, that $E^* \cup \{e\}$ is $(K,L)$-sparse for any edge $e \in E(V_{j} \cup V_{j+1})$. If $E^* \cap E(V_{j-1}\cup V_j)= \emptyset$, then it is easy to see, using inequality \eqref{conditionposet1 -k k l}, that $E^* \cup \{e\}$ is $(K,L)$-sparse for any edge $e \in E(V_{j} \cup V_{j+1})$. Hence, $\mathcal{I}_j$ is non-empty whenever $E(V_{j} \cup V_{j+1})$ is non-empty.

Let $E_1^{*}$ and $E_{2}^{*}$ be two maximally sparse subsets of $E(V_1 \cup \dots \cup V_{j})$, and let $\alpha_{j-1},\alpha'_{j-1} \in \mathcal{L}_{(K,L)}$ be such that $E_{\alpha_{j-1}} = E_1^*$, $E_{\alpha'_{j-1}} = E_2^*$. Suppose that $E\subseteq E(V_{j} \cup V_{j+1})$, and that $E \cup E_1^{*}$ is sparse. Our goal is to show that $E \cup E_{2}^{*}$ is also sparse.
We make a few remarks about $E \cup E_1^{*}$ first. Let $F\subseteq E$ and consider $V_{j}(F)$ and $V_{j+1}(F)$. We note that for any subset $A \subseteq E_1^{*}$, one has, by sparsity of $F \cup A$:
\begin{equation*}
|F \cup A|=|F| +|A| \leq  \Sigma_{i=1}^{j-1}k_i |V_i(A)| + k_{j} |V_{j}(A\cup F)| +  k_{j+1} |V_{j+1}(F)| -l_{{j},{j+1}}.
\end{equation*}
Hence, one has
\begin{equation*}
|F| \leq  \Sigma_{i=1}^{j-1} k_i |V_i(A)| + k_{j} |V_{j}(A\cup F)| +  k_{j+1} |V_{j+1}(F)| -l_{{j},{j+1}} - |A|.
\end{equation*}
Since $A$ is arbitrary, we see that:
\begin{equation*}
|F| \leq   k_{j+1} |V_{j+1}(F)| -l_{j,{j+1}} + \min_{A} \{ \Sigma_{i=1}^{j-1}k_i |V_i(A)| + k_{j} |V_{j}(A\cup F)| -|A|\}.
\end{equation*}
It is then clear that: 
\begin{equation*}
|F| \leq   k_{j+1} |V_{j+1}(F)| -l_{j,{j+1}} + \theta_{\alpha_{j-1}}(V_{j}(F)).
\end{equation*}
Now, for any subset $C\subseteq E_2^{*}$, we see that
\begin{align*}
|F \cup C|&=|F| + |C| \\
&\leq  k_{j+1} |V_{j+1}(F)| -l_{j,{j+1}} + \theta_{\alpha_{j-1}}(V_{j}(C\cup F)) + |C|\\
&=   k_{j+1} |V_{j+1}(F)| -l_{j,{j+1}} + \theta_{\alpha'_{j-1}}(V_{j}(C\cup F)) + |C|\\
&\leq   k_{j+1} |V_{j+1}(F)| -l_{j,{j+1}} + \Sigma_{i=1}^{j-1} k_i|V_i(C)| + k_{j}|V_{j}(C\cup F)|\\
&- |C| + |C|\\
&\leq  \Sigma_{i=1}^{j+1} k_{i}|V_{i}(C\cup F)|-l_{j,{j+1}} 
\end{align*}
We applied the observation that if $X\subseteq Y$ then $\theta_{\alpha_{j-1}}(X) \leq \theta_{\alpha_{j-1}}(Y)$ in the second line, with $X=V_{j}(F)$ and $Y = V_{j}(C\cup F)$. We have used Lemma \ref{surplus-independent} to go from the second line to the third. Since any subset of $E_2^{*} \cup E$ is of the form $F\cup C$, this shows that $E_2^{*} \cup E$ is sparse, completing the proof of the Lemma.
\end{proof}

\end{document}
